\documentclass[11pt]{article}

\usepackage[T1]{fontenc}
\usepackage[utf8]{inputenc}
\usepackage{lmodern}
\usepackage{geometry}
\geometry{margin=1in}
\usepackage{hyperref}
\usepackage{booktabs}
\usepackage{longtable}
\usepackage{array}
\usepackage{amsmath}
\usepackage{amssymb}
\usepackage{listings}
\usepackage{xcolor}
\setlength{\emergencystretch}{3em}

\hypersetup{
  colorlinks=true,
  linkcolor=blue!50!black,
  citecolor=blue!50!black,
  urlcolor=blue!50!black
}

\lstset{
  basicstyle=\ttfamily\small,
  breaklines=true,
  frame=single,
  columns=fullflexible
}

\newcommand{\code}[1]{\texttt{#1}}
\newcommand{\HE}{\textup{HE}}
\newcommand{\PeTTa}{\textup{PeTTa}}
\newcommand{\CeTTa}{\textup{CeTTa}}

\title{Semantic Contracts for \HE{} $\leftrightarrow$ \PeTTa{} Translation\\
\large Shared Core, Compatibility Profiles, and Production Boundaries}
\author{Zar Goertzel}
\date{May 2026}

\begin{document}
\maketitle

\begin{abstract}
This document states the semantic contract for translating between Hyperon
Experimental (\HE{}) MeTTa and \PeTTa{} MeTTa.  The central point is that both
systems use MeTTa syntax, but they do not always give the same operational
meaning to every surface form.  A faithful translator must therefore do more
than parse and reprint S-expressions: it must preserve observable behavior on a
shared core, lower source conveniences to explicit target forms where possible,
and classify helper or profile-specific surfaces explicitly when portability
requires more than syntax.

We describe the shared and non-shared language surfaces, the main translation
rules in both directions, the status of quote/eval/reduce, \code{test},
\code{hyperpose}, type queries, state, imports, and space operations, and the
validation obligations for production use.  The companion engineering
specification lives with the translator implementation at
\code{hyperon/translators/specs/he-petta-translation-semantics.md}.
\end{abstract}

\tableofcontents

\section{Problem}

\HE{} and \PeTTa{} are both MeTTa runtimes, but MeTTa syntax alone does not fix
all runtime behavior.  A file that ends in \code{.metta} may be:
\begin{enumerate}
\item a source program intended for default \PeTTa{};
\item an \HE{} program intended for a reference \HE{} runtime;
\item a generated \HE{}-facing artifact intended for \code{petta --he};
\item an artifact that supplies helper definitions to preserve source behavior;
\item an extension program that relies on host facilities outside core \HE{}.
\end{enumerate}

The translator's job is to keep these categories explicit.  Classifying every
generated artifact as pure \HE{} would be incorrect.  Classifying every
difference as an untranslatable failure would also be incorrect, because many
differences have principled lowerings.

\paragraph{Positive example.}
The \PeTTa{} source form
\begin{lstlisting}
!(test (length (collapse (demo-peano 300))) 301)
\end{lstlisting}
can be translated to an explicit \HE{}-facing computation:
\begin{lstlisting}
!(test (let $__tr_tuple_1 (collapse (demo-peano 300))
         (size-atom $__tr_tuple_1))
       301)
\end{lstlisting}
This is ordinary portable lowering.

\paragraph{Negative example.}
The \PeTTa{} source form
\begin{lstlisting}
!(test (quote (+ 1 2)) (+ 1 2))
\end{lstlisting}
depends on \PeTTa{}'s visible quote behavior.  A target runtime whose
\code{quote} returns a visible wrapper cannot be claimed to agree merely
because both sides parse the same text.

\section{Semantic Lanes}

The translator uses the following semantic lanes.

\begin{longtable}{>{\raggedright\arraybackslash}p{0.22\linewidth}
                  >{\raggedright\arraybackslash}p{0.46\linewidth}
                  >{\raggedright\arraybackslash}p{0.22\linewidth}}
\toprule
Lane & Meaning & Example \\
\midrule
\endhead
Shared core & Behavior is expected to agree on the target engines under
consideration. & \code{if}, \code{case}, many \code{let}, \code{match}, and
arithmetic examples. \\
Portable lowering & Source surface is rewritten to target forms without
\PeTTa{}-only assumptions. & \code{length (collapse e)} to \code{let} plus
\code{size-atom}. \\
Translator helper & Generated artifact needs a small helper definition to
preserve source behavior. & \code{test}, \code{assertEqualToEval},
\code{quoted-syntax}. \\
\PeTTa{} \HE{} compatibility & Intended for \code{./run.sh --he}; may rely on
profile behavior outside \HE{} core. & Callable-return application and selected
quote/eval/reduce compatibility. \\
\PeTTa{} extension & Useful runtime surface outside the strict portability
claim. & Python FFI, host resources, selected path/indexing support. \\
\bottomrule
\end{longtable}

This table is part of the correctness contract.  A helper-class artifact may be
correct for its intended target while remaining unsuitable as a backend-agnostic
\HE{} conformance example.

\section{\HE{} Surface}

The \HE{} target surface used by the translator is centered on explicit
evaluation and binding constructs.  Important target forms include:
\begin{center}
\begin{tabular}{ll}
\toprule
Construct & Role \\
\midrule
\code{chain} & Sequenced evaluation with result binding. \\
\code{collapse}, \code{collapse-bind} & Collection of nondeterministic results. \\
\code{superpose}, \code{superpose-bind} & Re-entry of alternatives. \\
\code{eval}, \code{evalc}, \code{unquote} & Explicit evaluation of syntax/data. \\
\code{case}, \code{switch} & Pattern-directed branching. \\
\code{size-atom} & Explicit tuple size/count surface. \\
\code{function}, \code{return} & Minimal function-return surface. \\
\code{metta} & Interpreter call with expected type and space. \\
\bottomrule
\end{tabular}
\end{center}

The translator should not turn implementation extensions into \HE{} core
requirements.  If a target runtime provides an extra surface, that surface must
be named as a profile or extension.

\section{\PeTTa{} Surface}

\PeTTa{} source programs use the same S-expression syntax plus Prolog-backed
runtime conveniences.  Some forms are shared.  Others require explicit
translation.

\begin{longtable}{>{\raggedright\arraybackslash}p{0.32\linewidth}
                  >{\raggedright\arraybackslash}p{0.58\linewidth}}
\toprule
\PeTTa{} surface & Translation status \\
\midrule
\endhead
\code{let}, \code{let*}, \code{if}, \code{case}, \code{match},
\code{collapse}, \code{superpose} & Usually shared or nearly identity. \\
\code{progn}, \code{prog1} & Lowered to sequenced \code{chain}/\code{let}
forms. \\
\code{foldall}, \PeTTa{}-style \code{foldl-atom} & Lowered to
\code{collapse} plus explicit accumulator/item variables. \\
\code{length (collapse e)} & Lowered to \code{let tuple (collapse e')
(size-atom tuple)}. \\
\code{test} & Preserved because its printed report is observable behavior. \\
\code{quote}, \code{eval}, \code{reduce} & Compatibility lane; see
Section~\ref{sec:quote}. \\
\code{hyperpose} & Lowered to \code{superpose} by default, preserved only with
explicit target intent. \\
Python, Git, and host FFI & Passed through or extension-classified; not pure
\HE{}. \\
\bottomrule
\end{longtable}

\section{\HE{} to \PeTTa{}}

The \HE{}-to-\PeTTa{} direction adapts explicit \HE{} sequencing and binding to
\PeTTa{} source forms:

\begin{center}
\begin{tabular}{lll}
\toprule
\HE{} input & \PeTTa{} output & Intent \\
\midrule
\code{(chain e \$x b)} & \code{(let \$x e b)} & Preserve sequencing. \\
\code{(collapse-bind e)} & \code{(collapse e')} & Preserve collection. \\
\code{(superpose-bind xs)} & \code{(superpose xs')} & Preserve alternatives. \\
\code{(nop e)} & \code{(let \$fresh e' ())} & Keep effect, discard result. \\
\code{(switch x bs)} & \code{(case x' bs')} & Translate branches. \\
\code{(function (return x))} & \code{x'} & Remove wrapper where safe. \\
\bottomrule
\end{tabular}
\end{center}

Fresh names use a reserved generated namespace such as \code{\$\_\_tr\_...}.
Correctness requires capture avoidance: generated names must not accidentally
bind or shadow source variables.

\section{\PeTTa{} to \HE{}}

The \PeTTa{}-to-\HE{} direction lowers source conveniences to explicit
\HE{}-facing forms:

\begin{center}
\begin{tabular}{>{\raggedright\arraybackslash}p{0.25\linewidth}
                >{\raggedright\arraybackslash}p{0.31\linewidth}
                >{\raggedright\arraybackslash}p{0.32\linewidth}}
\toprule
\PeTTa{} input & \HE{}-facing output & Intent \\
\midrule
\code{progn} & Nested sequencing & Preserve effects and final value. \\
\code{prog1} & Result capture plus sequencing & Preserve first result. \\
\code{foldall} & \code{collapse} plus \code{foldl-atom} & Make fold explicit. \\
\code{foldl-atom} & Explicit accumulator/item variables & Avoid implicit agg shape. \\
\code{length (collapse e)} & \code{let} plus \code{size-atom} & Portable lowering. \\
\code{test} & \code{test} & Preserve observable test surface. \\
\code{unique-atom (collapse e)} & \code{collapse (unique e')} & Use target uniqueness. \\
\code{hyperpose} & \code{superpose}, unless preserved & Drop scheduling, keep values. \\
\bottomrule
\end{tabular}
\end{center}

\section{Quote, Eval, Reduce}
\label{sec:quote}

Quote/eval/reduce are a common compatibility boundary.  In default \PeTTa{},
\begin{lstlisting}
!(quote (+ 1 2))
\end{lstlisting}
exposes the syntax value:
\begin{lstlisting}
(+ 1 2)
\end{lstlisting}
whereas some \HE{} implementations expose a visible quote wrapper:
\begin{lstlisting}
(quote (+ 1 2))
\end{lstlisting}

The \PeTTa{}-to-\HE{} translator preserves \PeTTa{}-visible behavior by
inserting a helper when needed:
\begin{lstlisting}
(= (quoted-syntax (quote $expr)) $expr)
\end{lstlisting}
and lowering:
\begin{lstlisting}
(quote e)   -> (quoted-syntax (quote e'))
(eval e)    -> (unquote (quote e'))
(reduce e)  -> (unquote (quote e'))
\end{lstlisting}

This is a compatibility translation, not a pure-\HE{} core claim.
\code{quoted-syntax} is a helper surface.  A backend-agnostic artifact must
avoid it, supply it explicitly, or replace it with a target-specific
quote/eval representation whose semantics are specified.

\section{Tests And Assertions}

\PeTTa{}'s \code{test} form prints an actual/expected report before returning
the assertion result.  That output is observable behavior.  Therefore the
translator preserves source \code{test} by default:

\begin{lstlisting}
!(test (+ 1 2) 3)
\end{lstlisting}

stays:

\begin{lstlisting}
!(test (+ 1 2) 3)
\end{lstlisting}

Quiet assertion helpers are appropriate for explicit profile tests.  Examples
include \code{assertEqualToEval}, \code{assertEqualToResult}, and
\code{assertAlphaEqualToResult}.  They are not a faithful replacement for a
source \code{test} whose printed report matters.

\section{Hyperpose}

\code{hyperpose} is \PeTTa{}'s parallel nondeterministic surface.  The default
portable translation lowers it to \code{superpose}, preserving the
nondeterministic values but not claiming to preserve parallel scheduling.

\begin{lstlisting}
(hyperpose xs) -> (superpose xs')
\end{lstlisting}

When the target runtime intentionally supports \code{hyperpose}, users may ask
the translator to preserve it:

\begin{lstlisting}
./translate.sh petta2he --preserve-hyperpose input.metta output.metta
\end{lstlisting}

Such a file is no longer the default pure-portable lowering.  It is a
runtime-profile artifact.

\section{Type Queries}

The translator depends on unknown type behavior being stable.  The intended
behavior is:

\begin{lstlisting}
!(get-type $x)      ; %Undefined%
!(get-metatype $x)  ; Variable
\end{lstlisting}

Returning a fresh type variable for \code{(get-type \$x)} is unsound for
translation: that variable can escape into stored type facts and later unify
with concrete types, producing branch-order-dependent routes.  Unknown types
should remain unknown; metatypes distinguish variables.

\section{Imports, Spaces, And State}

Translation over imports is supported in recursive and bundle modes.  Space and
state operations are part of the semantic boundary because they can expose
mutation order, duplicate atoms, backend indexing, and host resources.

\PeTTa{} named state and \HE{} native state share spellings such as
\code{new-state}, \code{get-state}, and \code{change-state!}, but they are not
the same portable contract.  A \PeTTa{} form such as:
\begin{lstlisting}
!(bind! counter (new-state 0))
\end{lstlisting}
uses \code{counter} as a mutable cell name.  \HE{} native state allocates an
explicit state handle.  The \PeTTa{}-to-\HE{} direction therefore lowers named
state through translator helpers instead of treating it as identity syntax.

The rule is simple: preserve or lower what has a specified target meaning, and
classify the rest.  Python handles, Git imports, MORK/PathMap handles, runtime
statistics, and other host resources are not pure \HE{} merely because the
syntax is readable by an \HE{} parser.

\section{Evidence And Gates}

There are several distinct evidence layers:

\begin{center}
\begin{tabular}{ll}
\toprule
Gate & Evidence \\
\midrule
\code{translate.sh --test} & Unit and real-file translator behavior. \\
\code{run\_he\_profile\_suite.sh} & PeTTa \code{--he} runtime profile behavior. \\
\code{run\_translator\_survey.sh} & Corpus correctness and performance budget. \\
\code{check\_generated\_he\_portability.sh} & Cross-engine classification. \\
Lean translator files & Proof-oriented common-fragment correspondence. \\
\bottomrule
\end{tabular}
\end{center}

These gates answer different questions.  A green \PeTTa{} \code{--he} survey is
not automatically a cross-engine proof.  A helper-surface classification is not
a runtime failure.  A proof over the common fragment does not license arbitrary
host extensions.

\section{Production Criteria}

A translation rule is production-ready only when:
\begin{enumerate}
\item its source and target semantic lanes are named;
\item positive and negative examples are recorded;
\item fresh variables are capture-avoiding;
\item generated helper names avoid source-symbol collisions;
\item observable output is preserved or explicitly classified;
\item helper surfaces are named as helpers, not \HE{} core;
\item translator tests pass;
\item representative generated artifacts run under their intended target;
\item cross-engine differences are classified when portability is claimed.
\end{enumerate}

Performance work must preserve these correctness obligations.  A faster
translation that changes source behavior is a regression.

\section{Conclusion}

The \HE{} $\leftrightarrow$ \PeTTa{} translator should be understood as a
semantic bridge between related MeTTa profiles, not as a string rewriter between
identical languages.  Its most important discipline is classification: shared
core behavior, portable lowering, helper surfaces, compatibility profiles, and
extensions must stay distinct.  With that discipline, the translator can serve
both engineering use and formalization work without making false portability
claims.

\end{document}
